\chapter{Graph Architecture and Node Specification}
\label{sec:graph}

This chapter describes the system as a graph: a set of nodes, each with one job,
connected by edges that show what flows between them. The goal is that this graph
can be built node-for-node in a framework such as LangGraph~\cite{langgraph}, with nothing left to
guess between what the thesis says and what the code does.

\section{Design Rationale}

Two design choices set this system apart from a standard single-pass RAG pipeline.

\textbf{Retrieval can repeat.} A standard pipeline searches once and answers once,
so it only works if that one search happens to find everything needed. Here, the
planner (N5) can call retrieval tools again and again, and the Sufficiency Gate
sends it back until the evidence is actually enough. Multi-hop reasoning is not a
prompting trick; it is simply how the loop is built.

\textbf{Only certain nodes can add evidence.} Nodes N6a--N7 are the only ones
allowed to bring anything into the evidence buffer (N8), and the generator (N9) can
only read from that buffer --- it has no other way to see the corpus. This
separation is what makes the Faithfulness Gate work: it only has to check claims
against one fixed source of truth, instead of guessing whether the generator made
something up.

\begin{figure}[p]
\centering
\resizebox{!}{0.72\textheight}{%
\begin{tikzpicture}[
  font=\sffamily,
  nd/.style={rectangle, rounded corners=3pt, draw, line width=0.7pt, minimum width=3.6cm, minimum height=1.15cm, align=center, inner sep=3pt, font=\sffamily\small},
  enc/.style={nd, fill=encblue, draw=encline},
  ret/.style={nd, fill=retbrown, draw=retline},
  fus/.style={nd, fill=fusteal, draw=fusline},
  svc/.style={nd, fill=svcorange, draw=svcline},
  sup/.style={nd, fill=supviolet, draw=supline},
  att/.style={nd, fill=attgreen, draw=attline},
  fin/.style={nd, fill=darkslate, draw=darkline},
  gate/.style={diamond, aspect=2.0, draw=gateline, line width=0.7pt, fill=gatered, align=center, inner sep=1pt, font=\sffamily\scriptsize},
  ar/.style={-{Latex[length=2mm]}, line width=0.6pt},
  loop/.style={-{Latex[length=2mm]}, line width=0.7pt, draw=gateline, dashed},
  back/.style={-{Latex[length=2mm]}, line width=0.6pt, draw=attline, dashed},
  lbl/.style={font=\sffamily\scriptsize, inner sep=2pt}
]

\node[enc] (n3) at (-4.6,0)    {\textbf{N3} Corpus Ingestion\\\scriptsize PDF $\rightarrow$ page rasters};
\node[enc] (n1) at (0,0)       {\textbf{N1} Query Intake\\\scriptsize Parse \& decompose};
\node[enc] (n2) at (4.6,0)     {\textbf{N2} Query Encoder\\\scriptsize Intent + filters};

\node[ret] (n4) at (-4.6,-2.3) {\textbf{N4} Dual Index Build\\\scriptsize ColPali + SPECTER2};
\node[sup] (n5) at (0,-2.3)    {\textbf{N5} Agentic Planner\\\scriptsize ReAct state machine};

\node[ret] (n6a) at (-4.6,-4.9){\textbf{N6a} TextRetriever\\\scriptsize Narrative chunks};
\node[ret] (n6b) at (0,-4.9)   {\textbf{N6b} TableRetriever\\\scriptsize Locate + cell parser};
\node[ret] (n6c) at (4.6,-4.9) {\textbf{N6c} FigureRetriever\\\scriptsize ColPali patches};

\node[ret] (n7) at (4.6,-6.9)  {\textbf{N7} Region Zoom\\\scriptsize Crop \& re-encode};

\node[fus] (n8) at (0,-9.0)    {\textbf{N8} Cross-Modal Linking\\\scriptsize Evidence buffer $\mathcal{B}_t$};

\node[gate] (g1) at (0,-10.9)  {\textbf{G1} Sufficiency\\Gate};

\node[svc] (n9) at (0,-12.7)   {\textbf{N9} Grounded Generator\\\scriptsize Answers from $\mathcal{B}_t$ only};

\node[gate] (g2) at (0,-14.5)  {\textbf{G2} Faithfulness\\Gate};

\node[att] (n10a) at (-4.6,-16.6){\textbf{N10a} Token Attrib.\\\scriptsize Integrated gradients};
\node[att] (n10b) at (0,-16.6)   {\textbf{N10b} Bounding Boxes\\\scriptsize Cell coords + patch sim.};
\node[att] (n10c) at (4.6,-16.6) {\textbf{N10c} DAG Trace\\\scriptsize Retrieval log};

\node[sup] (n11) at (0,-18.8)  {\textbf{N11} Report Synthesis\\\scriptsize Merge answer + proof};
\node[fin] (n12) at (0,-20.6)  {\textbf{N12} Final Output\\\scriptsize Verifiable answer};

\draw[ar] (n1) -- (n2);
\draw[ar] (n1) -- (n5);
\draw[ar] (n2.south) |- ([yshift=0.3cm]n5.east);
\draw[ar] (n3) -- (n4);

\begin{scope}[on background layer]
  \node[draw=retline, dashed, rounded corners=5pt, line width=0.6pt, fill=retbrown!25,
        fit=(n6a)(n6b)(n6c)(n7), inner sep=0.3cm] (tier) {};
\end{scope}
\node[lbl, anchor=south, text=retline, align=center] at (-2.3,0 |- tier.south) {retrieval tier:\\reads both N4 indexes};
\draw[ar] (n4.south) -- (n4.south |- tier.north) node[lbl, midway, left]{indexes};

\draw[ar] (n5.south) -- ([xshift=0.9cm]n6a.north);
\draw[ar] (n5) -- (n6b);
\draw[ar] (n5.south) -- ([xshift=-0.9cm]n6c.north);
\draw[ar] (n6c) -- (n7);

\draw[ar] (n6a.south) -- (-4.6,-7.9) -| (n8.north);
\draw[ar] (n7.south) -- (4.6,-7.9) -| (n8.north);
\draw[ar] (n6b) -- (n8);

\draw[ar] (n8) -- (g1);
\draw[ar] (g1) -- node[lbl, midway, right]{sufficient} (n9);
\draw[ar] (n9) -- (g2);

\draw[loop] (g1.east) -- node[lbl, pos=0.5, above, black]{insufficient} (7.3,-10.9) -- (7.3,-2.3) -- ([yshift=-0.3cm]n5.east);
\draw[loop] (g2.east) -- node[lbl, pos=0.5, above, black]{unfaithful} (7.3,-14.5) -- (7.3,-12.7) -- (n9.east);

\draw[ar] (g2.south) -- node[lbl, pos=0.55, right]{faithful} (0,-15.7);
\draw[ar] (0,-15.7) -- (-4.6,-15.7) -- (n10a.north);
\draw[ar] (0,-15.7) -- (4.6,-15.7) -- (n10c.north);
\draw[ar] (0,-15.7) -- (n10b.north);

\draw[back] (n10a.south) -- (-4.6,-18.0) -- (0,-18.0) -- (n11.north);
\draw[back] (n10c.south) -- (4.6,-18.0) -- (0,-18.0);
\draw[back] (n10b) -- (n11);
\draw[ar] (n11) -- (n12);

\end{tikzpicture}%
}
\caption{Full graph architecture of the proposed Agentic Multimodal RAG system.
Blue nodes perform ingestion and query encoding; brown nodes are the retrieval
tools that hold exclusive authority to introduce evidence, grouped in the shaded
retrieval tier because all of them query the two indexes built by N4; the teal node maintains
the evidence buffer that every downstream node reads from; the orange node is the
grounded generator; violet nodes are the planner and the synthesiser; green nodes
produce the three attribution artefacts. Red diamonds are the two validation gates
and red dashed edges are the two correction loops. Green dashed edges carry
attribution artefacts forward to synthesis. The planner may traverse the retrieval
tier arbitrarily many times before the Sufficiency Gate permits generation.}
\label{fig:graph}
\end{figure}

\section{Node-by-Node Specification}

Here is how one query moves through the graph, start to finish.

\begin{quote}
\textbf{Running example.} \emph{``Does the reported accuracy gain in Section~4 hold
for the low-resource subset, and which table reports it?''}
\end{quote}

This question needs two different kinds of evidence: a claim from the text and a
number from a table. A single search cannot answer it alone, which is exactly why
it is used here.

\subsection{Ingestion Path (N3--N4)}

This part runs once per corpus, offline, before any query arrives.

\paragraph{N3 --- Corpus Ingestion.}
Every PDF page is converted into an image, and its text is separately pulled out
and split into chunks. Pages become images rather than parsed layouts because
tables are exactly where OCR fails --- merged cells and multi-column headers break
text extraction. An image keeps the table's structure intact.

\paragraph{N4 --- Dual Index Build.}
Two search indexes are built from the same set of pages: one holds ColPali
visual embeddings of each page image (for finding tables and figures), the other
holds SPECTER2 text embeddings of each chunk, stored in FAISS~\cite{johnson2019faiss}
(for finding claims). Every retrieval tool (N6a--N6c, N7) reads from these two
indexes. Because both indexes point back
to the same pages, a match in one can lead directly to the matching spot in the
other --- which is what makes jumping from a text claim to its supporting table
possible.

\subsection{Query Encoding Path (N1--N2)}

\paragraph{N1 --- Query Intake.}
The question is broken into smaller steps before any search happens. For the
running example, that is two steps: find the accuracy claim, then find the table
that reports the low-resource breakdown. Doing this upfront gives the planner a
plan to follow instead of improvising one.

\paragraph{N2 --- Query Encoder.}
The query is embedded, and filters are pulled out of it --- here, ``Section 4'' and
``low-resource subset.'' This has to finish before any retrieval starts, since
every search the planner runs afterward depends on these filters.

\subsection{Planning and Retrieval (N5--N7)}

\paragraph{N5 --- Agentic Planner.}
This is the only node that decides what happens next. It keeps track of the plan,
what has already been searched, and what evidence has been found, then picks one
retrieval tool to run. After seeing the result, it either searches again or moves
on to the Sufficiency Gate. It never adds evidence itself --- only a retrieval tool
can do that.

\paragraph{N6a --- TextRetriever.}
Searches the text index for narrative chunks matching the query. On the running
example, this is the first search: it finds the sentence claiming the accuracy
gain, along with the page it appears on.

\paragraph{N6b --- TableRetriever.}
Finds a table in two steps: it searches the visual index for pages whose best
match is a table, then runs a table-structure recogniser (such as the Table
Transformer) on that region to recover rows, columns and the position of every
cell. Cell text is read from the PDF at those positions. It keeps track of exactly
where each cell sits on the page. This matters because the bounding boxes generated later (N10b) are only
possible if this spatial position is kept --- a plain-text version of the table
would lose it.

\paragraph{N6c --- FigureRetriever (ColPali).}
Searches the visual index to find charts, plots, or dense tables without needing
OCR. It also keeps the match scores it computed, because those scores are reused
later to draw bounding boxes.

\paragraph{N7 --- Region Zoom.}
Crops in on part of a page and re-encodes it at higher resolution. This exists
because finding the right page and reading a specific row inside it need different
levels of detail --- a page-level match tells you which page has the table, not
which row answers the question. On the running example, this runs after N6c finds
the right page, zooming into the low-resource row so its cells become readable.

\subsection{Evidence Buffer and Sufficiency Gate (N8, G1)}

\paragraph{N8 --- Cross-Modal Linking.}
Everything retrieved so far is converted into one shared format and added to the
evidence buffer, along with where it came from --- which page, which tool, which
step. Items that share a page, or where one refers to the other (a sentence citing
``Table~3'' and the table captioned ``Table~3''), are linked, so that the generator
receives connected evidence rather than a loose pile. This buffer is the only thing the generator will ever see; nothing after this
point looks at the corpus again.

\paragraph{G1 --- Sufficiency Gate.}
Before generation is allowed, a check runs: is there enough evidence here to
actually answer the question? On the running example, the first check fails --- the
text claim is there, but the table that would confirm or deny it for the
low-resource subset is not. Control goes back to N5, which fetches the table. Only
once both pieces are present does the gate let the process continue.

This check happens \emph{before} generation, not after. A standard pipeline would
hand the incomplete evidence to the generator and hope it notices something is
missing --- which is exactly how confident, wrong answers get produced.

\subsection{Generation and Faithfulness Gate (N9, G2)}

\paragraph{N9 --- Grounded Generator.}
An open-weight vision-language model (such as Qwen2-VL) run locally, so that N10a
can compute gradients through it. It writes the answer using only what is in the
evidence buffer. It cannot search on
its own and cannot fill in numbers from memory --- everything it says has to trace
back to something already retrieved.

\paragraph{G2 --- Faithfulness Gate.}
Checks every factual claim in the draft answer against the evidence buffer. If the
answer states something that is not actually backed by retrieved evidence, the
gate fails and N9 is sent back to rewrite that part using only what is really
there. Because the buffer is the only source allowed, each claim is checked
against one fixed, finite list of retrieved items, rather than against everything
the model might believe about the world.

\subsection{Attribution and Output (N10--N12)}

\paragraph{N10a --- Token Attribution.}
Highlights which retrieved sentences most influenced each part of the answer, using
integrated gradients computed through the N9 generator.

\paragraph{N10b --- Bounding Boxes.}
Draws a box --- $[x_{\min}, y_{\min}, x_{\max}, y_{\max}]$ --- around the exact spot
on the page a claim came from. For a table cell, the box is the cell position N6b
recorded; for a figure, it encloses the highest-scoring patches from N6c (or N7).
This only works because N6b kept cell positions and N6c kept its match scores.

\paragraph{N10c --- DAG Trace.}
Logs what actually happened: which tools ran, in what order, and how many times the
Sufficiency Gate sent things back. On the running example, this log shows two
search rounds and one rejection --- so the process can be checked, not just
trusted.

\paragraph{N11 --- Report Synthesis.}
Combines the answer with its three proof artefacts into one readable report.

\paragraph{N12 --- Final Output.}
What the user actually sees: the answer, the highlighted source text, a box on the
source page, and the log of how it was found.

\section{Control Loops}

There are exactly two loops in this graph, and each catches a different kind of
mistake.

\paragraph{The retrieval loop (not enough evidence).}
Triggered by the Sufficiency Gate. Sends control back to N5 to search again. This
stops the system from confidently answering using incomplete evidence.

\paragraph{The regenerate loop (wrong evidence used).}
Triggered by the Faithfulness Gate. Sends control back to N9 to rewrite. This stops
the system from stating something the evidence does not actually support.

Both loops have a limit on how many times they can repeat. If that limit is
reached, the system says the evidence is not enough instead of guessing. For a
literature review tool, this is the right trade-off: ``I could not confirm this''
is useful. A confident wrong answer is not.

\section{Mapping to Research Objectives}

Table~\ref{tab:node-mapping} shows which objective each part of the graph is
responsible for.

\begin{table}[htbp]
\centering
\small
\begin{tabular}{llp{5.4cm}}
\toprule
\textbf{Nodes} & \textbf{Objective} & \textbf{Contribution} \\
\midrule
N3, N4, N6c & Objective 2 & Vision-document retrieval over rasterised pages \\
N1, N2, N5, N7 & Objective 1 & Iterative agentic planning and multi-hop navigation \\
N8, Sufficiency Gate & Objective 3 & Evidence verification prior to generation \\
N9, Faithfulness Gate & Objective 3 & Claim-level grounding of generated output \\
N10a--N10c, N11 & Objective 4 & Spatial, textual, and procedural attribution \\
\bottomrule
\end{tabular}
\caption{Mapping of graph nodes to research objectives.}
\label{tab:node-mapping}
\end{table}